\phantomsection\addcontentsline{toc}{section}{Nanotechnology}
\ECchapter{27}{Nanotechnology}{Designing materials through scale effects}{ECGreen}

\ECObjectives{
\begin{itemize}
\item Explain why surface and quantum effects become dominant at nanometre scales.
\item Compare top-down and bottom-up fabrication using accuracy, throughput and defect criteria.
\item Quantify surface-area effects and define measurement, containment and lifecycle requirements.
\end{itemize}}

\begin{multicols}{2}
\begin{engineeringnews}
Below about one hundred nanometres, geometry and quantum physics can modify optical, electrical, chemical and
mechanical behaviour. Nanotechnology uses these scale effects in electronics, coatings, catalysts, sensors,
medicine and energy systems, but performance must be balanced against reproducibility and exposure risk.
\end{engineeringnews}

\begin{ecarticle}
A nanometre is one billionth of a metre. A nanoscale object may behave differently from the same material in bulk
form. For geometrically similar particles, surface area varies with the square of size while volume varies with
the cube. The surface-area-to-volume ratio therefore increases as characteristic length decreases. Interfaces,
adhesion, oxidation and surface reactions can then dominate the behaviour of the complete object.

For a spherical particle of diameter $d$,
\[
\frac{S}{V}=\frac{6}{d}.
\]
Halving the diameter doubles this ratio. This can increase catalytic activity or dissolution rate, but it may also
increase chemical reactivity, agglomeration and biological exposure.

A second scale effect is quantum confinement. When charge carriers are restricted to dimensions comparable with
their characteristic wavelengths, allowed energy levels change. Quantum dots can therefore emit different colours
according to size, and nanoscale conductors may display transport properties absent from larger components.

Top-down fabrication removes material from a larger structure using lithography, etching or precision machining.
It provides accurate patterns and compatibility with semiconductor processes, but cost and defect sensitivity rise
as dimensions decrease. Bottom-up fabrication assembles atoms, molecules or particles through synthesis,
deposition or self-organisation. It can form very small structures over large areas, although controlling size
distribution, orientation and contamination remains difficult.

Characterisation is part of the design loop. Conventional optical microscopy cannot resolve features far below the
wavelength of visible light, so electron microscopy, atomic-force microscopy and spectroscopy are used to measure
size, morphology, roughness and composition. A nominal particle diameter is insufficient if the batch contains a
wide distribution or large agglomerates.

Risk assessment must consider size, shape, surface chemistry, solubility and persistence rather than mass alone.
Responsible engineering therefore includes closed handling, local extraction, exposure monitoring, waste control,
recycling and end-of-life scenarios. A high-performance nanomaterial is not successful if its manufacture cannot be
controlled or verified.
\end{ecarticle}

\begin{tcolorbox}[title=\sffamily\bfseries Engineering insight,colback=yellow!10,colframe=ECOrange]
At the nanoscale, manufacturing tolerance can be of the same order as the designed feature. Metrology and process
capability are therefore as important as the nominal geometry.
\end{tcolorbox}

\begin{engineeringfocus}
\textbf{Scale, fabrication and verification}
\begin{center}
\resizebox{0.96\linewidth}{!}{%
\begin{tikzpicture}[font=\scriptsize,>=Latex,node distance=4mm]
\node[draw=ECBlue,fill=ECLightBlue,rounded corners] (req) {target property};
\node[draw=ECGreen,fill=ECGreen!10,rounded corners,right=of req] (fab) {nano-fabrication};
\node[draw=ECPurple,fill=ECPurple!10,rounded corners,right=of fab] (met) {metrology};
\node[draw=ECOrange,fill=ECOrange!10,rounded corners,below=of fab] (risk) {exposure and lifecycle};
\draw[->,thick] (req)--(fab); \draw[->,thick] (fab)--(met);
\draw[->,thick] (met) to[bend left=25] node[above]{feedback} (req);
\draw[->,thick,dashed] (fab)--(risk); \draw[->,thick,dashed] (risk)--(req);
\end{tikzpicture}}
\end{center}
\end{engineeringfocus}

\begin{keyfigures}
\begin{tabularx}{\linewidth}{@{}Xr@{}}
1 nanometre & $10^{-9}$ m\\
Typical nanoparticle & 1--100 nm\\
Top-down & material removal\\
Bottom-up & assembly\\
Critical controls & metrology and exposure
\end{tabularx}
\end{keyfigures}

\begin{tcolorbox}[title=\sffamily\bfseries Surface area increases as size decreases,colback=white,colframe=ECGreen]
\begin{center}
\begin{tikzpicture}
\begin{axis}[width=0.94\linewidth,height=45mm,xlabel={Particle diameter (nm)},ylabel={Relative surface area},
grid=major,ticklabel style={font=\scriptsize},label style={font=\scriptsize},x dir=reverse]
\addplot[thick,mark=*] table[x=diameter,y=area,col sep=comma]{data/EC27/surface.csv};
\end{axis}
\end{tikzpicture}
\end{center}
\footnotesize The curve illustrates a geometric scaling law; real performance also depends on accessible surface and agglomeration.
\end{tcolorbox}

\begin{applicationexercise}{Particle size and accessible surface}
A coating is available with spherical particles of diameter $40\,\mathrm{nm}$ or $200\,\mathrm{nm}$.
\begin{enumerate}
\item Calculate $S/V=6/d$ for both sizes in $\mathrm{m^{-1}}$.
\item Determine the ratio between the two values.
\item If 30\% of the surface of the smaller particles is lost through agglomeration, calculate the effective gain.
\item Explain one performance benefit and one exposure risk associated with the smaller particles.
\end{enumerate}
\end{applicationexercise}

\begin{vocabulary}
\begin{tabularx}{\linewidth}{@{}lX@{}}
nanoscale & échelle nanométrique\\
nanoparticle & nanoparticule\\
thin film & couche mince\\
lithography & lithographie\\
etching & gravure\\
self-assembly & auto-assemblage\\
quantum dot & boîte quantique\\
agglomeration & agglomération\\
coating & revêtement\\
process capability & capabilité du procédé
\end{tabularx}
\end{vocabulary}

\begin{questions}
\item Why does the surface-area-to-volume ratio rise as size decreases?
\item Compare one advantage and one limitation of top-down and bottom-up fabrication.
\item Why are electron or atomic-force microscopes often required?
\item Why can agglomeration reduce the expected performance of nanoparticles?
\item Why is particle mass alone insufficient for evaluating health risk?
\end{questions}

\begin{engineeringchallenge}
Specify a nanostructured anti-corrosion coating for an aluminium component. Define the target property, particle
size range, fabrication route, two measurable acceptance criteria, one sampling plan, two exposure controls and an
end-of-life strategy. Include a short risk--benefit table and justify whether the nanoscale solution is necessary.
\end{engineeringchallenge}

\begin{speaking}
Do the potential benefits of nanoparticles justify deployment before all long-term environmental effects are
known? Discuss evidence, precaution, regulation and lifecycle responsibility.
\end{speaking}
\end{multicols}

\subsection*{Sources}
\ECsource{European Commission -- Nanomaterials}{https://environment.ec.europa.eu/topics/chemicals/nanomaterials_en}
\ECsource{National Nanotechnology Initiative}{https://www.nano.gov/}